\documentclass[a4paper,11pt]{article}
        \usepackage[top=15mm,bottom=18mm,left=16mm,right=16mm]{geometry}
        \usepackage{fontspec}
        \usepackage{xeCJK}
        \setmainfont{Times New Roman}
        \setCJKmainfont{Yu Gothic}
        \setmonofont{Consolas}
        \setCJKmonofont{Yu Gothic}
        \usepackage{booktabs}
        \usepackage{longtable}
        \usepackage{tabularx}
        \usepackage{array}
        \usepackage{enumitem}
        \usepackage{hyperref}
        \usepackage{listings}
        \usepackage{xcolor}
        \hypersetup{
          colorlinks=true,
          linkcolor=blue,
          urlcolor=blue,
          pdftitle={車体物理・電気モデル本体},
          pdfauthor={Codex}
        }
        \lstset{
          basicstyle=\ttfamily\small,
          breaklines=true,
          columns=fullflexible,
          keepspaces=true,
          frame=single,
          framerule=0.2pt,
          rulecolor=\color{black},
          showstringspaces=false
        }
        \setlength{\parskip}{0.42em}
        \setlength{\parindent}{1em}
        \renewcommand{\arraystretch}{1.15}
        \title{12. 車体物理・電気モデル本体}
        \author{solar\_ws0129-main}
        \date{2026-07-29}
        \begin{document}
        \maketitle

        \section*{対象}
        \begin{tabularx}{\linewidth}{>{\raggedright\arraybackslash}p{0.18\linewidth} X}
        \toprule
        項目 & 内容 \\
        \midrule
        ファイル & \path|mpc_solarcar/model.py| \\
        種別 & Python \\
        区分 & model \\
        役割 & 空力、転がり、坂、PV、MPPT、drive/regen 効率、battery IV、SoC 更新を統合した vehicle model。 \\
        起動文脈 & planner と simulation の数理コア。 \\
        \bottomrule
        \end{tabularx}

        \section*{このファイルを読む理由}
        空力、転がり、坂、PV、MPPT、drive/regen 効率、battery IV、SoC 更新を統合した vehicle model。

        \begin{itemize}[leftmargin=1.6em]
        \item SolarCarModel と Params が中心。
\item electrical\_balance が planner 側で最も多く呼ばれる。
\item resistive\_forces、battery\_iv、soc\_step が各所から再利用される。
        \end{itemize}

        \section*{呼び出し関係}
        \subsection*{どこから呼ばれるか}
        \begin{itemize}[leftmargin=1.6em]
        \item \path|mpc_node.py|
\item \path|solar_state_node.py|
\item \path|scripts/solar_sim.py|
\item \path|estimator.py|
        \end{itemize}

        \subsection*{次に読むと理解が進むファイル}
        \begin{itemize}[leftmargin=1.6em]
        \item \path|mpc_solarcar/utils_maps.py|
        \end{itemize}

        \subsection*{同一リポジトリ内の主要依存}
        \begin{itemize}[leftmargin=1.6em]
        \item \path|mpc_solarcar/utils_maps.py|
        \end{itemize}

        \section*{ソース冒頭の把握}
        モジュール docstring または先頭意図の要約:

        モジュール docstring は明示されていない。

        \subsection*{代表コード断片}
        \begin{lstlisting}[language=Python]
)

@dataclass
class Params:
    dt: float=600.0
    rho: float=1.18
    air_density_mode: str='constant'
    air_density_reference_pressure_pa: float=101325.0
    CdA: float=0.13
    Crr: float=0.002
    Crr_per_wheel: float=0.0
    m: float=250.0
    g: float=9.80665
    P_aux: float=60.0
    P_aux_stopped: float=60.0
    P_aux_night: float=0.0
    aux_night_ghi_threshold_wm2: float=20.0
    gear_eta: float=0.98
    gear_ratio: float=6.0
    wheel_radius: float=0.28
    wheel_count: int=4
    driven_wheel_count: int=2
\end{lstlisting}


        \section*{主要構造の読み取り}
        主要クラスは Params, SolarCarModel。 主要関数は eff\_drive, eff\_regen, select\_drive\_mode, R\_int, pv\_balance, pv\_power\_mppt, charge\_efficiency, soc\_step。

        \subsection*{主要クラス・関数}
            \begin{longtable}{p{0.07\linewidth} p{0.16\linewidth} p{0.25\linewidth} p{0.40\linewidth}}
            \toprule
            行 & 種別 & 名前 & 補足 \\
            \midrule
            17 & class & \path|Params| &  \\
73 & class & \path|SolarCarModel| &  \\
10 & function & \path|_is_symbolic| & args: x \\
74 & function & \path|__init__| & args: self, drive\_map\_path, regen\_map\_path, Rint\_map\_path, params, panel\_eff\_map\_path, mppt\_eff\_map\_path, drive\_map\_eco\_path, drive\_map\_power\_path, regen\_map\_eco\_path, regen\_map\_power\_path, ocv\_soc\_map\_path \\
124 & function & \path|eff_drive| & args: self, v\_ms, tau\_nm \\
137 & function & \path|eff_regen| & args: self, v\_ms, tau\_nm \\
150 & function & \path|_update_mode_limits| & args: self \\
158 & function & \path|_select_mode| & args: self, v\_ms, tau\_nm \\
169 & function & \path|select_drive_mode| & args: self, v\_ms, tau\_nm \\
172 & function & \path|R_int| & args: self, T\_C, z \\
179 & function & \path|pv_balance| & args: self, G\_poa, T\_cell\_C \\
222 & function & \path|pv_power_mppt| & args: self, G\_poa, T\_cell\_C \\
225 & function & \path|_scaled_slope_pct| & args: self, slope\_pct \\
228 & function & \path|charge_efficiency| & args: self, P\_pack \\
235 & function & \path|soc_step| & args: self, z, P\_pack, dt\_sec \\
262 & function & \path|ocv_from_soc| & args: self, z \\
271 & function & \path|load_ocv_map| & args: self, path \\
280 & function & \path|air_density| & args: self, ambient\_temp\_c, elevation\_m \\
297 & function & \path|resistive_forces| & args: self, v\_ms, slope\_pct, headwind\_ms \\
344 & function & \path|battery_iv| & args: self, P\_pack, z, Tbat\_C \\
368 & function & \path|mech_power| & args: self, v\_ms, slope\_pct, headwind\_ms, inertial\_power\_w \\
389 & function & \path|auxiliary_power| & args: self, aux\_power\_w \\
396 & function & \path|scheduled_auxiliary_power| & args: self \\
405 & function & \path|torque_from_mech| & args: self, P\_mech, v\_ms, wheel\_radius, ratio \\
416 & function & \path|electrical_balance| & args: self, v\_ms, slope\_pct, z, Tbat\_C, G\_poa, Tcell\_C, headwind\_ms, aux\_power\_w, inertial\_power\_w, ambient\_temp\_c, elevation\_m \\
            \bottomrule
            \end{longtable}

        \subsection*{ファイルを上から読んだときの定義順}
            \begin{longtable}{p{0.07\linewidth} p{0.84\linewidth}}
            \toprule
            行 & 内容 \\
            \midrule
            3 & 例外処理を伴う try ブロックを実行する。 \\
10 & 関数 \_is\_symbolic を定義する。 \\
17 & クラス Params を定義する。 \\
73 & クラス SolarCarModel を定義する。 \\
            \bottomrule
            \end{longtable}

        \subsection*{import 群の詳細}
            \begin{longtable}{p{0.07\linewidth} p{0.33\linewidth} p{0.50\linewidth}}
            \toprule
            行 & import 文 & このファイルで読む理由 \\
            \midrule
            1 & \path|import math| & clip、sqrt、sin/cos、有限判定など数値ロジックに使うため。 このファイル内での主な使用位置は L329, L331, L336, L362。 \\
2 & \path|import numpy as np| & 数値配列、clip、補間、統計計算を行うため。 このファイル内での主な使用位置は L135, L148, L266, L269, L292, L347, L447。 \\
4 & \path|import casadi as ca| & casadi モジュールを利用するため。 このファイル内での主な使用位置は L6, L11, L12, L126, L128, L130, L139, L143, ...。 \\
7 & \path|from dataclasses import dataclass| & dataclasses から dataclass を読み込み、このファイルの処理を組み立てるため。 このファイル内での主な使用位置は L16。 \\
8 & \path|from .utils_maps import read_eff_map, read_Rint_map, read_map, bilinear_interp, read_1d_map| & 効率マップ・抵抗マップ読込補間 から read\_eff\_map, read\_Rint\_map, read\_map, bilinear\_interp, read\_1d\_map を読み込み、このファイルの内部処理を分担させるため。 実体ファイルは mpc\_solarcar/utils\_maps.py。 このファイル内での主な使用位置は L81, L82, L93, L95, L97, L99, L101, L106, ...。 \\
            \bottomrule
            \end{longtable}

        \section*{関数・クラスを上から順に解説}
        \subsection*{L10 関数 \path|_is_symbolic|}
\begin{itemize}[leftmargin=1.6em]
  \item 定義: \path|_is_symbolic(x)|
  \item docstring: 明示されていない。
  \item このブロックが直接呼ぶ主な関数・メソッド: hasattr, is\_symbolic, isinstance
  \item 戻り値の要点: ca is not None and (isinstance(x, (ca.SX, ca.MX)) or (hasattr(x, 'is\_symbolic') and x.is\_symbolic()))
\end{itemize}
\begin{enumerate}[leftmargin=1.8em]
\item ca is not None and (isinstance(x, (ca.SX, ca.MX)) or (hasattr(x, 'is\_symbolic') and x.is\_symbolic())) を返す。
\end{enumerate}

\subsection*{L17 クラス \path|Params|}
            \begin{itemize}[leftmargin=1.6em]
              \item 定義: \path|Params(bases=なし)|
              \item docstring: 明示されていない。
              \item このブロックが直接呼ぶ主な関数・メソッド: 特に明示されていない。
              \item 戻り値の要点: 戻り値が無いか、return 文が明示されていない。
            \end{itemize}
            \begin{enumerate}[leftmargin=1.8em]
            \item dt に 600.0 を代入する。
\item rho に 1.18 を代入する。
\item air\_density\_mode に 'constant' を代入する。
\item air\_density\_reference\_pressure\_pa に 101325.0 を代入する。
\item CdA に 0.13 を代入する。
\item Crr に 0.002 を代入する。
\item Crr\_per\_wheel に 0.0 を代入する。
\item m に 250.0 を代入する。
\item g に 9.80665 を代入する。
\item P\_aux に 60.0 を代入する。
\item P\_aux\_stopped に 60.0 を代入する。
\item P\_aux\_night に 0.0 を代入する。
            \end{enumerate}

\subsection*{L73 クラス \path|SolarCarModel|}
            \begin{itemize}[leftmargin=1.6em]
              \item 定義: \path|SolarCarModel(bases=なし)|
              \item docstring: 明示されていない。
              \item このブロックが直接呼ぶ主な関数・メソッド: Params, R\_int, \_is\_symbolic, \_scaled\_slope\_pct, \_select\_mode, \_update\_mode\_limits, abs, air\_density, atan, auxiliary\_power, battery\_iv, bilinear\_interp
              \item 戻り値の要点: 戻り値が無いか、return 文が明示されていない。
            \end{itemize}
            \begin{enumerate}[leftmargin=1.8em]
            \item 関数 \_\_init\_\_ を定義する。
\item 関数 eff\_drive を定義する。
\item 関数 eff\_regen を定義する。
\item 関数 \_update\_mode\_limits を定義する。
\item 関数 \_select\_mode を定義する。
\item 関数 select\_drive\_mode を定義する。
\item 関数 R\_int を定義する。
\item 関数 pv\_balance を定義する。
\item 関数 pv\_power\_mppt を定義する。
\item 関数 \_scaled\_slope\_pct を定義する。
\item 関数 charge\_efficiency を定義する。
\item 関数 soc\_step を定義する。
            \end{enumerate}











        \section*{処理の流れ}
        \begin{enumerate}[leftmargin=1.7em]
        \item ソース中の主要関数を通じて処理を進める。
        \end{enumerate}

        \section*{このファイルの位置づけ}
        この資料群では、\path|mpc_solarcar/model.py| を
        \textbf{model}の中核として扱う。
        したがって、ここで扱う処理は単独で完結せず、
        前段の profile / launch / telemetry と後段の planner / logger / report へ繋がる。

        \end{document}
